\documentclass[11pt]{article}

% ===============================
% Packages
% ===============================
\usepackage[margin=1in]{geometry}
\usepackage{amsmath, amsthm, amssymb, amsfonts}
\usepackage{mathtools}
\usepackage{enumitem}
\usepackage{hyperref}
\usepackage{graphicx}
\usepackage{physics}
\usepackage{bm}
\usepackage{microtype}
\usepackage{parskip}
\usepackage{float}
\hypersetup{
    colorlinks=true,
    linkcolor=blue,
    urlcolor=blue,
    citecolor=blue
}

% ===============================
% Custom Commands
% ===============================
\newcommand{\R}{\mathbb{R}}
\newcommand{\proj}{\mathrm{proj}}
\newcommand{\cone}{\mathcal{C}}
\newcommand{\Hil}{d_{\mathrm{H}}}

% ===============================
% Title
% ===============================
\title{\textbf{Note on Non-linear Perron-Frobenius dynamics in Gradient Flows via Hilbert metric}}
\author{Xinyang (Elizabeth) Wen}
\date{November 29, 2025}

% ===============================
% Document
% ===============================
\begin{document}

\maketitle


% =======================================================
% 1. Introduction / Motivation
% =======================================================
\section{Motivation}
% -------------------------------------------------------
% Your content here.
% Introduce the phenomenon, motivations, background.
% -------------------------------------------------------

My central motivation is the following:
in many common learning problems, the data distribution naturally
concentrates on a low-dimensional manifold,
so naturally the gradient descent iteration can be viewed as a form of
diffusion on this underlying specific geometric structure, depending on data. Once a loss
functional is introduced, the diffusion acquires a preferred direction,
and a partial order arises naturally through the loss function(Lyapunov structure).

In diffusion, changing the geometry and changing the generator often produce equivalent effects on diffusion behavior (dual operations).So I initially guess it might
behaves like a non-linear Perron-Frobenius
system (at least in expectation).

The purpose of this note is to record empirical observations—from Swiss Roll to MNIST
(and SVM in the future to avoid the difficulties on MNIST)

\section{Observation on Swiss Roll}
\subsection{model:}
In this section I summarize the empirical behavior observed on the
Swiss Roll regression/classification experiments.  The task is to learn
a smooth target function on the Swiss Roll manifold using gradient
descent on a simple single-layer model with a nonnegative
parametrization (linear map followed by an activation), trained with
various loss functions.  The Swiss Roll is first embedded into a
high-dimensional ambient space (up to dimension $4096$ in some runs), in
order to examine the effect of dimensionality and noise on the projective
geometry of the parameter trajectory.

Across all Swiss Roll experiments, the parameter trajectory exhibits a
surprisingly consistent projective structure.  The main observations are
as follows.
Formulas:
\[
\hat{y} = \sigma(w^\top x),
\qquad w \in \mathbb{R}^d,
\]
where $\sigma$ is chosen so that $\sigma'(z)\ge 0$ for all $z$
(e.g.\ ReLU, softplus), ensuring that the parameter dynamics remain in a
positive cone.
Specifically, in my experiments, the $\sigma$ is Softplus activation function : \[
\text{Softplus}(x) = \ln\left(1 + e^x\right)
\]
Given a loss function $L(\hat{y}, y)$ (either squared loss or Huber
loss), the gradient of the objective for a single data point is
\[
\nabla_w L = L'(\hat{y},y)\,\sigma'(w^\top x)\,x.
\]
Thus each update takes the form
\[
w_{t+1}
= w_t - \eta\, L'(\sigma(w_t^\top x), y)\,\sigma'(w_t^\top x)\,x,
\]
and for a batch $\{(x_i,y_i)\}$,
\[
w_{t+1}= w_t - \eta \sum_i
L'(\sigma(w_t^\top x_i), y_i)\,\sigma'(w_t^\top x_i)\,x_i.
\]
No kernelization or feature engineering is used; the only source of
nonlinearity is the activation $\sigma$ and the choice of loss.  The
positivity of $\sigma'$ ensures that the trajectory $\{w_t\}$ evolves
inside a positive cone, which is the geometric setting relevant for
Hilbert metric and Perron--Frobenius--type behavior.

\subsection{Square loss(MSE)}
\subparagraph{\textbf{Result}:}
\begin{enumerate}
    \item \textbf{full batch GD}:Under small noise, with the squared loss, the Hilbert distance to the final vector decreases almost geometrically to float 0 distance to the final vector immediately after several initial steps (alignment).
    \item \textbf{Alternating cone contraction}: Under specific condition(in my toy model it's a range depending on random seeds, embedding dimension, and Noise), the trajectory behaves like a 2-cycle contraction(The problem is I don't have a tool to analysis the period larger than 2).
    \item \textbf{Stochastic GD}: Under Stochastic GD, the iterations converges almost linearly (in the plot, Hilbert distance to the final vector). It behaves as a Perron-Frobenius Operator under this experiment (in the meaning of float numbers).
\end{enumerate}

\subparagraph{\textbf{Analysis(Guessing):}}

\begin{enumerate}
    \item I think the Softplus makes the whole matrix positive strictly and the Hessian matrix at the local minimum point is semi-definite. So in our model it's approximate to be a linear operator.
    \item  It's like a period-2 contraction. Once their Birkhoff contraction ratio is closed, the leading contraction vector alternates. I verify this by tracking the largest eigenvalue during training. The data shows the largest eigenvalues decrease at 2 cycles.
    \item I didn't find a good explanation for this so that this can also explains some interesting phenomenon in MNIST models(I haven't finished  my notes yet). The problem is this only works on square loss. When I switched to Huber loss, the stochastic GD cannot guarantee this. Another weird thing is on MNIST with neuron network with hidden layers model, the full-batch can guarantee non-expansion while this disappears under stochastic GD.
\end{enumerate}

\subsection{Huber loss.(Full-Batch GD)}
Huber loss
\[
L_\delta(a)=
\begin{cases}
\frac{1}{2}a^2, & |a|\le\delta,\\[3pt]
\delta(|a|-\frac{\delta}{2}), & |a|>\delta,
\end{cases} \quad  a = y - \hat{y}.
\]
\subparagraph{\textbf{Result:}}
\begin{enumerate}
    \item \textbf{full batch GD}: \ The Hilbert trajectory shows a two-stage pattern under small noise after 3-5 initial iteration: an initial
almost-arithmetic decrease, followed by near-geometric contraction.  This mirrors the linear-to-quadratic transition in the loss itself.
    \item \textbf{Alternating cone contraction}:The periodic-like trajectory appears more on the Huber loss models, even I increase the batch size to 4096.
    \item \textbf{Stochastic GD}: The Stochastic GD significantly smooth the alternating cone contraction, but it's almost impossible to see a non-expansive phenomenon on Hilbert metric.
\end{enumerate}
\subparagraph{\textbf{Analysis:}}
\begin{enumerate}
    \item \textbf{Linear}: after several initial vectors, since the difference between the predict y is not very closed to the exact y. Under Huber loss, the gradient becomes sign-like → update behaves like a piecewise positive, non-smooth cone map with no global Hessian.

    \textbf{MSE}:when the loss is closed, according to the formula of Huber loss, the loss function is closed to the MSE loss.
    \item I think it's because the Huber loss is piecewise, so the two parts has different eigenvector, as a result it switches cones frequently during the training.(with similar contraction ration I guess?)
    \item \textbf{more smooth} The stochastic GD introduces more noise so it might weaken the deterministic of the 2-cycle.

    \textbf{Fail restore non-expansion} I think it's because the noise is not enough to restores the difference between the two cycles led by the formula of Huber loss
\end{enumerate}
\section{Current observations on MNIST models }
\textbf{I haven't organized this part very well yet }

For MNIST, I use a standard $L$-layer fully-connected network with ReLU
activations and cross-entropy loss (regularization and other settings follow
standard practice).  The goal is again to monitor the last-layer trajectory in
the Hilbert metric.

\paragraph{(1) Stepwise behavior.}
Both full-batch and stochastic gradient descent show a characteristic
``plateau--drop''  pattern: long regions where the Hilbert distance stays almost
constant, followed by sudden drops over a short steps.

\paragraph{(2) Effect of depth.}
Increasing network depth makes the drops larger and the plateaus longer.
In full-batch training this effect is especially clear.

\paragraph{Summary}
These observations are preliminary; I am still organizing the experiments and
will provide a more complete analysis later.

\end{document}
